\documentclass[10pt,twocolumn]{article}
\usepackage[margin=0.6in,top=0.55in,bottom=0.55in]{geometry}
\usepackage{amsmath,amssymb}
\usepackage{booktabs}
\usepackage{hyperref}
\usepackage{enumitem}
\usepackage{microtype}
\usepackage{titlesec}
\usepackage{url}
\usepackage{setspace}
\setstretch{0.9}

\titlespacing*{\section}{0pt}{4pt}{1.5pt}
\titlespacing*{\subsection}{0pt}{2.5pt}{0.5pt}
\titlespacing*{\paragraph}{0pt}{2pt}{0.6em}
\renewcommand{\arraystretch}{0.9}
\setlength{\columnsep}{14pt}
\setlength{\abovecaptionskip}{3pt}
\setlength{\belowcaptionskip}{0pt}
\setlength{\textfloatsep}{6pt}
\setlength{\intextsep}{6pt}
\setlength{\parskip}{0pt}
\abovedisplayskip=4pt \belowdisplayskip=4pt
\abovedisplayshortskip=2pt \belowdisplayshortskip=2pt

\hypersetup{colorlinks=true,linkcolor=blue,urlcolor=blue,citecolor=blue}

\title{\Large\textbf{TicTacPro: Exactly Solving a Combinatorial Piece-Placement Game\\
and Stress-Testing Deep Q-Learning Against It}}

\author{\normalsize Abdul Khan \quad
\small\texttt{abdulskhan47@gmail.com} \quad
\small\url{https://github.com/abdulllkhan/ticTACpro-sim}}
\date{\vspace{-1.2em}\small July 2026}

\begin{document}
\maketitle
\thispagestyle{empty}

%% ---------------------------------------------
\begin{abstract}
\noindent
We study TicTacPro, a combinatorial board game extending Tic-Tac-Toe with
three piece sizes per player and two winning conditions. An exhaustive
alpha-beta solver proves the game is a \emph{first-player win from every one
of the 27 opening moves} (700M nodes, 246\,s on 20 cores). Notably,
time-limited heuristic minimax self-play had first converged to the
\emph{opposite} conclusion --- a deterministic second-player win --- a
spurious result the exact solver refutes, illustrating how agreement among
strong but budget-bounded agents can masquerade as a game-theoretic proof.
We further document a block-priority bug that recurred across nine
independent code paths, and train a Dueling Double DQN (792K parameters)
with a curriculum of rule-based opponents, reaching 100\% win rates against
random and BullseyeAgent opponents. A post-hoc audit surfaces
schedule-calibration and checkpoint-selection pitfalls, reported as
negative results.
\end{abstract}

%% ---------------------------------------------
\section{Game Description}

TicTacPro is played on a $3\!\times\!3$ board. Each player holds nine pieces:
three Small (S), three Medium (M), and three Large (L). On each turn a player
places one piece on any cell. A cell has one slot per size, so a cell can
hold up to three pieces. A player wins by achieving either:
\begin{itemize}[noitemsep,topsep=2pt,leftmargin=*]
  \item \textbf{3-in-a-row (same size):} three pieces of identical size along
        any row, column, or diagonal; or
  \item \textbf{Bullseye:} one S, one M, and one L piece placed in the same
        cell.
\end{itemize}
A game lasts at most 18 plies, with up to 27 moves per ply --- substantially
harder to search than Tic-Tac-Toe, yet still small enough to solve exactly,
which proved essential (Section~\ref{sec:gametheory}).

%% ---------------------------------------------
\section{Game-Theoretic Analysis}
\label{sec:gametheory}

\subsection{Exact Solution: a First-Player Win}
\label{sec:exact}

We solved the game with a pure negamax alpha-beta search over the complete
game tree: no heuristic evaluation, no depth cap, terminal-only scoring
(win/loss/draw), with a transposition table, parallelised over the 27
opening moves on 20 CPU cores. \textbf{Every one of the 27 opening moves is a
proven win for the first player (RED).} Individual openings required 9.2M
(TR-M) to 53.7M (BC-S) node evaluations; 246\,s wall-clock total. The
per-opening table and solver are in the repository
(\texttt{exact\_solve.py}, \texttt{exact\_solve\_results.json}).

\subsection{The Search-Horizon Trap}
\label{sec:horizon}

Before the exact solve, we ran time-limited iterative-deepening minimax
(3\,s/move, heuristic leaf evaluation, late-move reductions, aspiration
windows) against itself. Five independent games produced an \emph{identical}
16-move sequence in which BLUE (the second player) wins via a ``sequential
bullseye fork'': build two sizes of a bullseye in one cell, force RED to
spend a piece blocking, pivot to a fresh cell, repeat until RED runs out of
blocks. Its determinism made this look like a solution --- we initially
recorded the game as a second-player win.

The exact solver shows why this was wrong. Proving RED's forced win from the
minimax engine's own book opening (TC-L) requires $\sim$31.8M node
evaluations; the online engine evaluates only $\sim$150--300K nodes per move
--- two orders of magnitude short. The fork is a real and potent strategy
\emph{against bounded search}, but it is not the game value. Nor does
reactive play execute it safely: a rule-based BullseyeAgent implementing the
fork beats the rule-based anti-diagonal OptimalAgent 20/20 in \emph{both}
colour roles (heuristic dominance, not a second-player advantage), yet loses
10/10 to the same bounded minimax via double attacks no reactive policy can
prevent. \textbf{Lesson:} deterministic agreement among strong, budget-bounded
agents is weak evidence about game value; exhaustive search reversed the
answer.

%% ---------------------------------------------
\section{The Block-Priority Bug Family}

Nine move-selection code paths shared a subtle correctness bug: a single
\texttt{block\_mv} variable was set on the \emph{first} blocking move found
in scan order (Small$\to$Large, cell 0$\to$8). When both a \emph{bullseye
threat} (opponent has 2/3 sizes in one cell) and a \emph{line threat}
(opponent has 2/3 of a same-size row) coexist, scan order --- not danger ---
determined which was blocked.

With flat board index $\mathrm{idx} = 9r + 3c + \mathrm{sz}$, two partner
indices $i_1,i_2$ of a win line lie in the same cell iff
$\lfloor i_1/3\rfloor = \lfloor i_2/3\rfloor$, identifying bullseye threats
in $O(1)$. The fix tracks the two threat classes in separate variables
(\texttt{bullseye\_block}, \texttt{line\_block}) and prefers the bullseye
block. It was propagated to all nine affected paths: OptimalAgent,
NeuralMCTS fast-path, NeuralMCTS move sorting, the C move-sorting extension,
MCTS move sorting, the Python rollout policy, the C rollout, the C
single-move selector, and the ParallelMCTS fast-path. Each site received a
regression test. The recurrence of one latent flaw across nine independent
implementations of ``win $>$ block $>$ rest'' is itself the finding: shared
decision logic should be shared code.

%% ---------------------------------------------
\section{Deep Q-Learning Agent}

The 61-dimensional state tensor is normalised to the current player's
perspective: channels 0--26 encode the current player's pieces (one per
cell-size slot), 27--53 the opponent's, 54--59 the two players' remaining
piece counts (scaled by $1/3$), and channel 60 flags whether the current
player is RED (the first mover).

\begin{table}[h]
\centering\small
\setlength{\tabcolsep}{4pt}
\begin{tabular}{ll}
\toprule
\textbf{Component} & \textbf{Value} \\
\midrule
Architecture & Dueling Double DQN \\
Parameters & 792,604 \\
Action space & 27 (9 cells $\times$ 3 sizes) \\
State dimension & 61 \\
Precision & BF16 AMP + \texttt{torch.compile} \\
Replay buffer & 10M (GPU-side PER) \\
Symmetry augmentation & 8-fold ($D_4$), at collection \\
Parallel collectors & 16 CPU workers (spawn) \\
Episodes (run 8, 2\,h) & 1,003,520 \\
Hardware & NVIDIA GB10 DGX Spark \\
\bottomrule
\end{tabular}
\caption{Training configuration for the final run.}
\label{tab:config}
\end{table}

\paragraph{Negamax Double-DQN target.} The zero-sum structure lets us write
the target from the opponent's perspective, with the online network selecting
among \emph{legal} actions and the target network evaluating:
\begin{equation*}
\hat{Q} = r - (1-d)\,\gamma\;
Q_{\theta^-}\!\bigl(s',\, \arg\max_{a' \in \mathcal{A}(s')} Q_{\theta}(s',a')\bigr),
\end{equation*}
where $r \in \{+1,0,-1\}$ is the terminal reward, $d$ the done flag, and
$\mathcal{A}(s')$ the legal-action mask. The opponent's best continuation is
\emph{subtracted}, matching the minimax Bellman equation.

\paragraph{8-fold symmetry.} The dihedral group $D_4$ yields 8 board
symmetries. Each collected transition is expanded into its 8 symmetric
variants \emph{before insertion} into the replay buffer (states via inverse
permutation gather, actions via forward permutation; count channels and the
first-mover flag are invariant), multiplying stored experience by $8\times$.

\paragraph{Systems optimisations.} A C rollout extension runs entire greedy
rollouts in one call ($4.3\times$ over Python); collector workers use a
pure-NumPy inference path ($19\times$ over PyTorch CPU); PER sampling runs
on-GPU via prefix-sum and \texttt{searchsorted}.

\paragraph{Curriculum.} Pure self-play never encounters the bullseye-fork
strategy, so the DQN never learns to block it. The final run therefore
replaced one player with the rule-based BullseyeAgent in a majority
($\approx$56\%, measured) of collected episodes, in both colour roles.

%% ---------------------------------------------
\section{Results}

\begin{table}[h]
\centering\footnotesize
\setlength{\tabcolsep}{3pt}
\begin{tabular}{lccc}
\toprule
\textbf{Match} & \textbf{W} & \textbf{D} & \textbf{L} \\
\midrule
DQN vs Random (both roles, 40g) & 40 & 0 & 0 \\
DQN vs Bullseye (both roles, 40g) & 40 & 0 & 0 \\
DQN (RED) vs Optimal (BLUE), 20g & 0 & 20 & 0 \\
DQN (BLUE) vs Optimal (RED), 20g & 0 & 0 & 20 \\
\bottomrule
\end{tabular}
\caption{Post-training tournament (rule-based agents).}
\label{tab:tournament}
\end{table}

Overall win rates: BullseyeAgent and DQN 66.7\% each, OptimalAgent 50.0\%,
Random 0.0\%. Three observations. \textbf{(1)~Curriculum transfer:} before
the BullseyeAgent curriculum the DQN lost every benchmarked game to it
(0/20); after, it wins 40/40 --- it learned to block the sequential bullseye
fork. \textbf{(2)~Colour asymmetry:} the DQN draws OptimalAgent when playing
first but loses 20/20 when playing second; it never saw the anti-diagonal
opening plan in training. \textbf{(3)~Tree search:} in the tournament series
with search agents, pure MCTS with win-then-block greedy rollouts reached a
91.7\% overall score (tournament v8), and replacing its random rollouts with
deterministic OptimalAgent rollouts made it \emph{worse}: rollout diversity
mattered more than rollout strength.

%% ---------------------------------------------
\section{Limitations (Post-Hoc Training Audit)}
\label{sec:limitations}

Auditing the pipeline after the final run surfaced three flaws.
\textbf{Stale-throughput schedules:} epsilon, PER-$\beta$, and LR schedules
assumed a hard-coded 120K gradient steps/hour; the run achieved $\sim$31K/h,
ending at $\varepsilon = 0.457$ and $\beta = 0.649$ --- no exploitation
phase, incomplete importance-sampling correction. Schedules should track
measured progress. \textbf{Saturating checkpoint metric:} the
best-checkpoint criterion (win rate vs.\ a deterministic heuristic) hit
100\% eleven minutes in and could never improve again, so the benchmarked
checkpoint reflects 143,872 of the 1M episodes; deterministic-vs-deterministic
evaluation also carries only 1 bit of signal per colour role.
\textbf{Silent loss signal:} in heuristic-opponent games only DQN-side
transitions are stored; when the opponent's move ends the game the DQN's
final transition keeps $r=0$, $d=0$, so losses must be learned through
bootstrapped opponent-to-move states --- exactly the states such games never
store.

%% ---------------------------------------------
\section{Conclusions}

Our agent combines the standard value-based stack --- DQN~\cite{mnih2015},
Double Q-learning~\cite{vanhasselt2016}, dueling heads~\cite{wang2016},
prioritized replay~\cite{schaul2016} --- with UCT~\cite{kocsis2006} baselines
and a classical alpha-beta solver; the need to verify agent consensus
exhaustively echoes solved-game efforts such as checkers~\cite{schaeffer2007}.
Contributions:

\begin{enumerate}[noitemsep,topsep=2pt,leftmargin=*]
  \item \textbf{Solved game:} TicTacPro is a first-player win from all 27
        openings; bounded minimax self-play had confidently produced the
        opposite answer.
  \item \textbf{Bug pattern:} one latent block-priority flaw recurred in nine
        independent implementations; dual-tracking threat classes is the
        reusable fix.
  \item \textbf{Curriculum beats self-play blindness:} a rule-based adversary
        in collection taught the DQN a defence self-play never discovered
        (0/20 $\to$ 40/40).
  \item \textbf{Audit your pipeline:} the flaws in
        Section~\ref{sec:limitations} were invisible in headline win rates
        and surfaced only through post-hoc log analysis.
\end{enumerate}

\paragraph{Reproducibility.} \texttt{exact\_solve.py} reproduces the
game-value proof ($\sim$4 min, 20 cores); \texttt{train\_spark.py} and
\texttt{tournament.py} in the repository reproduce the training run and
Table~\ref{tab:tournament}.

\begin{thebibliography}{9}\scriptsize
\setlength{\itemsep}{0pt}\setlength{\parsep}{0pt}\setlength{\topsep}{1pt}
\bibitem{mnih2015} V.~Mnih et al. Human-level control through deep
reinforcement learning. \emph{Nature}, 518:529--533, 2015.
\bibitem{vanhasselt2016} H.~van Hasselt, A.~Guez, D.~Silver. Deep
reinforcement learning with Double Q-learning. \emph{AAAI}, 2016.
\bibitem{wang2016} Z.~Wang et al. Dueling network architectures for deep
reinforcement learning. \emph{ICML}, 2016.
\bibitem{schaul2016} T.~Schaul, J.~Quan, I.~Antonoglou, D.~Silver.
Prioritized experience replay. \emph{ICLR}, 2016.
\bibitem{kocsis2006} L.~Kocsis, C.~Szepesv\'ari. Bandit based Monte-Carlo
planning. \emph{ECML}, 2006.
\bibitem{schaeffer2007} J.~Schaeffer et al. Checkers is solved.
\emph{Science}, 317:1518--1522, 2007.
\end{thebibliography}

\end{document}
